\section{Retained inputs, applications and attribution}\label{retained-inputs-applications-and-attribution}

The proofs use the elementary subgroup, quotient and isomorphism
theorems. Square recognition also uses the Euclidean algorithm and unique
factorization in $\Z[i]$, and the elementary point stabilizers in the
square reflection tiling described in its proof. The further classical
inputs are explicitly limited:
Burnside's basis theorem for finite \(p\)-groups; their nilpotence;
generation of the lower-central graded Lie ring by degree one and its
multilinear commutator rules; and the precise amalgam normal-form
theorem stated in the amalgam section. The Frattini, forest, finite-quotient and
torsion deductions are proved in full above. No finite group
classification, atlas census or numerical search result supplies a
premise to these theorems.

For applications to abstract regular polytopes we retain the following
standard correspondence {[}MS02, Section 2E{]}. A group generated by an
ordered list of nonidentity involutions, with the string commutation
rules and IP, is the full automorphism group of a regular abstract
polytope with that distinguished marking. Its rank is the length of the
list, its type consists of the adjacent product orders, and its number
of flags is the group order. A generator-preserving epimorphism between
such groups induces a rank- and adjacency-preserving covering of their
associated polytopes. These are geometric input theorems, not new
conclusions of the present paper. By contrast, the coset-geometric
recognition and coordinate-kernel Helly statements needed for the
hypertope branch are proved in the incidence section with their full hypotheses.

The supported-extension construction supplies a common proof interface
for the two-power constructions studied in {[}HFL20, HFL19a, HFL19b,
CFHS25{]}. The presentation-recognition theorem supplies the additional normality and
exact-base conditions needed to identify a compatible product with a complete
universal presentation; local construction alone does not do this.
The forest quotient gives the canonical group underlying the main
quotient and cover conclusions of {[}HS26{]}, and the finite-kernel
parameter space gives the exhaustive part of the order-1024
reconstruction associated with {[}GLP18{]}. Their exact source
statements, literal corrections, numerical certificates and named-object
identifications belong to the separate applications.

The supported-kernel induction, forest quotient and affine models are
reused complete proofs. The normal-kernel recognition theorem replaces
repeated finite-order comparisons, and the mix criterion separates joint
coordinate detection from generation of a fiber product. Every application
still proves its local defining relations, normality and base identifications.
Original source results and borrowed geometric theorems keep their attribution.

Existing Lean certificates cover declared components of the kernel,
supported-extension and Helly arguments. They do not presently formalize
this whole paper: in particular, the full forest quotient proof, the
affine model and the geometric recognition/application interfaces are
not all connected in Lean. Their written proofs above remain part of the
mathematical obligation. The accompanying certificate scope
distinguishes inherited checked declarations from any freshly checked
application components.


The mixed-prime application also retains Burnside's $p^a q^b$ solvability
theorem, the exact Lucchini cyclic-core bound as quoted with its hypotheses,
and the complete SmallGroups catalogue at order 384. The affine bridge
uses actual Reidemeister--Schreier and Tietze computations of two kernels,
not their abelianizations. The graph-symmetry application retains Frucht's
theorem and the complete SmallGroups catalogues at orders 40,48,56 for its
isolated rank-three minimum. These are application inputs; none is silently
used to prove the general lattice-coset or normal-kernel theorems.

All eleven complete required program/data files are retained once in this
volume: four files for the two-group classification (including its complete
named-presentation data and runner), six for the mixed-prime affine families,
and one for the graph-symmetry minimum. The three complete appendices print
those bytes. Operational drivers and all previously scoped Lean files are
separately inventoried. Existing successful semantic and computational
reviews are reused for unchanged bytes; this consolidation does not claim
to have reexecuted every old census or to have formally verified the new
arbitrary-parabolic lattice criterion.
